% Brutus --- headnote
% brutus, 46 BC, Rome (dedicated to M. Iunius Brutus)

\textit{The \textit{Brutus} is Cicero's history of Roman
oratory: the first attempt by any Roman to set down, in order
and with judgment, the whole succession of public speakers from
the obscure beginnings of the Republic to his own day. Cast as
a dialogue and dedicated to M. Iunius Brutus, it was written
early in 46 BC, in the enforced literary leisure of Caesar's
dictatorship --- after Pharsalus, in the months around the
death of Cato at Utica, when the Forum that had been Cicero's
whole life had fallen silent. Its occasion is grief: the death
the year before of Q. Hortensius Hortalus, the great rival of
Cicero's youth and the only orator who had ever stood beside
him at the head of the bar. The loss of Hortensius becomes the
loss of an art, and the loss of the art becomes the loss of the
free commonwealth in which alone, Cicero argues, eloquence can
live.}

\textit{The scene is Cicero's own garden, by a statue of Plato;
the speakers are Cicero, Brutus, and T. Pomponius Atticus, whose
recently published \textit{Liber Annalis} supplies the
chronological spine of the survey. After a Greek prelude (the
rise of eloquence at Athens, the line down to Demosthenes, and
the decline that followed him), Cicero turns to Rome and unrolls
a catalogue of some two hundred orators --- from the half-legendary
figures of the early Republic and the first named speaker, M.
Cornelius Cethegus, through Cato the Elder, Galba, the Gracchi,
and Lepidus Porcina (the first to give Latin prose a Greek finish),
to the central pairing of the dialogue: M. Antonius and L.
Licinius Crassus, the two masters of Cicero's boyhood, set side
by side and weighed against each other. Along the way Cicero
defends his governing thesis --- that the verdict of the expert
and the verdict of the crowd, rightly understood, come to the
same thing --- and pauses over the living, including a famous
tribute to the pure Latinity of Caesar himself.}

\textit{The history bends, at its close, into autobiography.
Cicero narrates his own formation --- his training in philosophy,
law, and rhetoric, his journey to Greece and Asia, his study
under Molon of Rhodes --- and his rise, after Hortensius, to the
first place at the bar. This personal turn carries the polemical
heart of the work: a defence of Cicero's own copious, fully
armed manner against the young men who called themselves
``Attici'' and prized a bare, thinned-down plainness --- C.
Licinius Calvus among them, and, by sympathy, Brutus too. The
dialogue ends not in triumph but in mourning. With the courts
closed and free speech extinguished under one man's power, Cicero
addresses Brutus directly, grieving for a young orator whose
course has been cut across by ``the wretched fortune of the
commonwealth,'' and the surviving text breaks off, fittingly, on
a broken sentence --- eloquence itself falling silent in the very
act of recording its own history.}
